\documentclass[11pt,a4paper]{article}
\usepackage[T2A]{fontenc}
\usepackage[utf8]{inputenc}
\usepackage[russian,english]{babel}
\usepackage{geometry,xcolor,array,booktabs,longtable,tabularx,tikz,fancyhdr,titlesec,enumitem,hyperref,lastpage,microtype}
\usepackage[sfdefault]{noto}
\usepackage[most]{tcolorbox}
\usepackage{pdflscape}
\usetikzlibrary{arrows.meta,calc,positioning,shapes.geometric,babel}

\geometry{left=18mm,right=18mm,top=21mm,bottom=20mm,headheight=20pt}
\setlength{\parindent}{0pt}
\setlength{\parskip}{5pt}
\setlength{\emergencystretch}{2em}
\setlist{nosep,leftmargin=*}
\definecolor{limblue}{HTML}{102A43}
\definecolor{limcyan}{HTML}{00A6A6}
\definecolor{limaccent}{HTML}{FF6B4A}
\definecolor{limmint}{HTML}{DFF7F5}
\definecolor{limlight}{HTML}{F1F5F8}
\definecolor{limpaper}{HTML}{FAFCFD}
\definecolor{limgray}{HTML}{627D98}
\hypersetup{colorlinks=true,linkcolor=limcyan,urlcolor=limcyan,pdftitle={LIM M2 --- Архитектурная спецификация процессора}}
\pagestyle{fancy}
\fancyhf{}
\fancyhead[L]{\colorbox{limblue}{\textcolor{white}{\strut\hspace{2mm}\textbf{LIM M2}\hspace{2mm}}}}
\fancyhead[R]{\small\textcolor{limgray}{PROCESSOR REFERENCE MANUAL / REVISION 0.2}}
\fancyfoot[L]{\small LIM M2 Reference Manual --- Редакция 0.2 (Октябрь 2026)}
\fancyfoot[R]{\small \thepage/\pageref{LastPage}}
\renewcommand{\headrulewidth}{0pt}
\titleformat{\section}{\Large\bfseries\color{limblue}}
  {\colorbox{limcyan}{\textcolor{white}{\strut\thesection}}}{0.75em}{}
\titleformat{\subsection}{\large\bfseries\color{limblue}}{\thesubsection}{0.7em}{}
\titlespacing*{\section}{0pt}{16pt}{8pt}
\titlespacing*{\subsection}{0pt}{5pt}{1pt}

\newcommand{\TBD}{\textcolor{limcyan}{\textbf{TBD}}}
\newcommand{\reg}[1]{\texttt{#1}}
\newcommand{\bits}[1]{\texttt{#1}}
\newcommand{\field}[1]{\textsf{#1}}
\newcommand{\opcodebox}[1]{\tikz[baseline=(n.base)]\node[fill=limblue,text=white,rounded corners=2pt,inner xsep=7pt,inner ysep=3pt,font=\ttfamily\bfseries] (n) {#1};}
\newtcolorbox{limnote}[1]{
  enhanced,breakable,colback=limlight,colframe=limcyan,boxrule=0pt,
  borderline west={3pt}{0pt}{limcyan},arc=1mm,left=4mm,right=4mm,
  top=3mm,bottom=3mm,title={#1},fonttitle=\bfseries\color{limblue},
  coltitle=limblue,attach title to upper={\par\smallskip}}
\newtcolorbox{specbox}{
  enhanced,colback=limpaper,colframe=limlight,boxrule=0.8pt,arc=2mm,
  left=3mm,right=3mm,top=2mm,bottom=2mm}

% Шаблон страницы инструкции: одна команда на отдельную страницу
\newcommand{\InstructionPage}[8]{%
  \clearpage
  \phantomsection\addcontentsline{toc}{subsection}{#1 --- #4}
  \begin{tcolorbox}[enhanced,colback=limblue,colframe=limblue,arc=3mm,
    left=5mm,right=5mm,top=4mm,bottom=4mm]
  \begin{minipage}[t]{0.68\textwidth}
    {\fontsize{30}{34}\selectfont\bfseries\color{white}#1}\par\vspace{1mm}
    {\large\color{white}#4}
  \end{minipage}\hfill
  \begin{minipage}[t]{0.25\textwidth}\raggedleft
    \tikz[baseline=(n.base)]\node[fill=limcyan,text=white,rounded corners=2pt,
      inner xsep=8pt,inner ysep=4pt,font=\ttfamily\bfseries] (n) {#2};
    \par\smallskip{\small\color{white}#3}
  \end{minipage}
  \end{tcolorbox}
  \vspace{3mm}
  \begin{specbox}
  \begin{tabularx}{\dimexpr\textwidth-8mm\relax}{>{\bfseries\color{limblue}}p{31mm}X}
    \toprule
    Мнемоника & \texttt{#1} \\
    Opcode & \bits{#2} (старшие 5 бит кода команды) \\
    Формат команды & 5 бит opcode, 3 бита \field{cfg} и аргументы переменного формата (у специализированных команд --- операнды сразу после opcode) \\
    Класс операции & #3 \\
    Статус фиксации & Базовый функционал утверждён; битовые поля аргументов и расширенные модификаторы находятся на стадии финализации \\
    \bottomrule
  \end{tabularx}
  \end{specbox}
  \subsection*{Назначение команды} #4.
  \subsection*{Архитектурная семантика}
  \begin{tcolorbox}[colback=limmint,colframe=limmint,arc=2mm,boxrule=0pt,left=4mm]
    \ttfamily #5
  \end{tcolorbox}
  \subsection*{Конфигурация и операнды} #6
  \subsection*{Влияние на регистр флагов (FREG)} #7
  \subsection*{Микрооперации исполнительного тракта}
  \begin{enumerate}[topsep=1pt,itemsep=0pt,parsep=0pt]
    \item Коммутация мультиплексоров тракта данных для выборки исходных операндов.
    \item Исполнение операции в специализированном функциональном блоке (АЛУ / память / переходы).
    \item Направление результата через арбитр записи (MUX 2) к целевому регистру-приёмнику.
    \item Синхронное обновление аппаратных флагов \reg{FREG} и служебных регистров (\reg{PTREG}, \reg{SREG}).
  \end{enumerate}
  Точная сетка микротактов и управляющие стробы: \TBD.
  \subsection*{Особенности реализации и примечания} #8
  \vfill{\footnotesize\color{limgray}Параметры с отметкой TBD соответствуют расширенным опциям спецификации.}
}

\begin{document}
\begin{titlepage}
  \pagecolor{limpaper}\color{limblue}
  \begin{tikzpicture}[remember picture,overlay]
    \fill[limblue] (current page.north west) rectangle
      ([xshift=58mm]current page.south west);
    \fill[limcyan] ([xshift=58mm]current page.north west) rectangle
      ([xshift=63mm]current page.south west);
    \draw[limlight,line width=0.7pt] ([xshift=77mm,yshift=-37mm]current page.north west)
      -- ([xshift=-16mm,yshift=-37mm]current page.north east);
  \end{tikzpicture}
  \vspace*{18mm}
  \hspace*{-7mm}\begin{minipage}{50mm}\centering
    {\color{white}\bfseries\large HIGH-PERFORMANCE}\par\smallskip
    {\color{limcyan}\bfseries\small 16-BIT ARCHITECTURE}
  \end{minipage}\par
  \vspace{31mm}
  \hspace*{68mm}\begin{minipage}{100mm}\raggedright
    {\fontsize{38}{42}\selectfont\bfseries LIM\par M2}\par
    \vspace{7mm}
    {\Large\bfseries\color{limcyan}Архитектурное руководство по процессору}\par
    \vspace{9mm}
    {\large Высокодетерминированная 16-разрядная процессорная архитектура с чистым мультиплексорным трактом данных и расширенной 24-битной адресацией.}\par
    \vspace{18mm}
    \begin{tikzpicture}[node distance=5mm]
      \foreach \x/\n in {0/R0,1.25/R1,2.5/R2,3.75/R3,5/R4,6.25/R5,7.5/R6,8.75/R7}
        \node[draw=limblue,rounded corners=1.5pt,minimum width=10mm,minimum height=8mm,
          fill=white,font=\bfseries\scriptsize] (r\n) at (\x,0) {\n};
      \draw[limcyan,line width=2pt,-{Latex[length=2mm]}] (0,-0.8) -- (8.75,-0.8);
      \draw[limblue,line width=1.2pt,-{Latex[length=2mm]}] (8.75,-1.25) -- (0,-1.25);
      \node[anchor=west,font=\small\bfseries,text=limgray] at (0,-1.85)
        {POINT-TO-POINT DATA ROUTING \enspace·\enspace ZERO BUS CONTENTION};
    \end{tikzpicture}
  \end{minipage}
  \vfill
  \hspace*{68mm}\begin{minipage}{100mm}\raggedright
    {\bfseries\color{limaccent}OFFICIAL REFERENCE SPECIFICATION}\par\smallskip
    {\large Редакция 0.2 --- Октябрь 2026 г.}\par\medskip
    {\small\color{limgray}Спецификация ядра LIM M2. Все права защищены.}
  \end{minipage}
  \vspace*{10mm}
\end{titlepage}
\nopagecolor

\tableofcontents
\clearpage

\section{Статус документа и архитектурные цели}
Настоящий документ является официальным техническим описанием и архитектурным стандартом микропроцессора \textbf{LIM M2}. 
Архитектура спроектирована с прицелом на бескомпромиссную надежность, математическую строгость и высокую тактовую частоту.

В основу тракта данных LIM M2 заложен фундаментальный принцип \textbf{полного отказа от внутренних шин третьего состояния (Tri-State / Hi-Z)}. Внутренние потоки данных построены исключительно на детерминированных двухтактных мультиплексорных деревьях (Point-to-Point MUX Routing), что гарантирует полное отсутствие сквозных токов, шинных конфликтов (bus contention) и паразитных емкостных задержек.

\begin{limnote}{Архитектурная нумерация команд}
Система команд LIM M2 использует непрерывную 5-битную сетку первичных кодов: 
\texttt{NOP}=0, \texttt{ADD}=1, \ldots, \texttt{MSP}=26. Коды 27--31 зарезервированы для аппаратных расширений следующего поколения.
\end{limnote}

\subsection{Условные обозначения и сокращения}
\begin{tabularx}{\textwidth}{>{\ttfamily}p{28mm}X}
\toprule
Обозначение & Архитектурный смысл \\
\midrule
dst & Целевой регистр-приёмник результата операции \\
src, srcA, srcB & Регистры-источники входных операндов \\
imm & Непосредственное скалярное значение, встроенное в тело инструкции \\
addr & Эффективный адрес памяти (16-битный базовый или 24-битный расширенный) \\
cfg & 3-битное поле функционального модификатора операции \\
ZF, CF, OF & Аппаратные флаги нуля, переноса и знакового переполнения регистра \reg{FREG} \\
\bottomrule
\end{tabularx}

\subsection{Единая философия формата инструкций}
Архитектура процессора LIM M2 строго придерживается универсальной парадигмы: \textbf{сначала всегда указывается целевой приёмник (куда), а затем входные операнды (откуда)}. Это обеспечивает строгий, логичный и интуитивно понятный синтаксис:
\begin{itemize}
  \item \textbf{Арифметические операции (3 операнда):} \texttt{ADD Rd, Rs1, Rs2} означает $\text{Rd} \leftarrow \text{Rs1} + \text{Rs2}$ (например, \texttt{ADD R0, R1, R2} вычисляет $R0 = R1 + R2$); аналогично для \texttt{SUB}, \texttt{MUL}, \texttt{DIV}.
  \item \textbf{Логические операции (2 операнда):} \texttt{AND Rd, Rs} означает $\text{Rd} \leftarrow \text{Rd} \ \text{AND}\ \text{Rs}$ (например, \texttt{AND R0, R1} вычисляет $R0 = R0 \ \& \ R1$); аналогично для \texttt{OR}, \texttt{NAND}, \texttt{NOR}, \texttt{XOR}, \texttt{XNOR}.
  \item \textbf{Побитовая инверсия NOT (1 операнд, 1-байтовая инструкция):} \texttt{NOT Rd} означает $\text{Rd} \leftarrow \sim\text{Rd}$ (например, \texttt{NOT R0} вычисляет $R0 = \sim R0$). Кодируется в 1 байте (\texttt{[7:3]} опкод, \texttt{[2:0]} регистр Rd), второй байт из памяти не запрашивается.
  \item \textbf{Межрегистровые пересылки:} \texttt{LRR Rd, Rs} означает прямое копирование из точки Б (источник \reg{Rs}) в точку А (приёмник \reg{Rd}): $\text{Rd} \leftarrow \text{Rs}$ (например, \texttt{LRR R0, R1} означает $R0 = R1$).
  \item \textbf{Загрузка из памяти (LMR):} первый аргумент \reg{Rd} задаёт целевой регистр процессора, куда сохраняются считанные данные.
\end{itemize}

\section{Архитектурный обзор процессора LIM M2}
\textbf{LIM M2} представляет собой высокоэффективный 16-разрядный микропрограммный процессор общего назначения с развитой целочисленной арифметикой. Ключевые архитектурные достоинства процессора включают:

\begin{itemize}
  \item \textbf{Чистая мультиплексорная маршрутизация данных:} Внутренний тракт полностью избавлен от третьего состояния. Передача данных осуществляется активными драйверами через каскад селекторов \textbf{MUX 1}, \textbf{DEST-DMUX} и \textbf{MUX 2}, что делает архитектуру идеально совместимой с современными полупроводниковыми техпроцессами (ASIC) и быстрыми ПЛИС (FPGA).
  \item \textbf{Двухаккумуляторный исполнительный узел (Dual Accumulator Staging):} Регистры \reg{ACCA} и \reg{ACCB} обеспечивают надёжную фиксацию операндов перед АЛУ, изолируя переходные процессы вычислений от основного регистрового файла.
  \item \textbf{Аппаратное каскадирование длинной арифметики (ALU Auto-Chaining):} Внутренний тракт АЛУ оснащён остаточным аккумулятором \reg{ALU\_RESIDUAL}. При возникновении переноса ($+1$ при сложении $>2^{16}-1$), заёма ($-1$ при разности $<0$) или избыточных старших разрядов умножения/деления, остаток фиксируется аппаратно и автоматически суммируется или вычитается при следующей математической операции. Это обеспечивает непрерывную многобайтовую арифметику (32, 48, 64+ бит) без отдельных команд ADC/SBB. Сброс цепочки выполняется командой \texttt{NOP}.
  \item \textbf{Прямой тракт пересылки (Reg-to-Reg Bypass):} Обходной канал от MUX 1 напрямую в MUX 2 позволяет осуществлять копирование данных между регистрами без задействования и блокировки вычислительных ресурсов АЛУ.
  \item \textbf{Расширенная 24-битная адресация памяти (до 16 МБ):} Специализированные 24-битные регистры-указатели \reg{PTREG} (указатель данных с автоинкрементом) и \reg{SREG} (указатель стека) позволяют процессору линейно адресовать обширные массивы памяти и стек до 16~Мбайт за пределами классического 64-килобайтного окна.
  \item \textbf{Высокочастотный потенциал:} Отсутствие емкостной нагрузки от распределённых шин позволяет ядру достигать базовой рабочей частоты \textbf{200~МГц}, а при трансляции в специализированный кремний (ASIC) ориентироваться на частоты \textbf{700~МГц и выше}.
\end{itemize}

\clearpage
\begin{landscape}
\thispagestyle{empty}
\subsection{Концептуальная схема тракта данных}
\begin{center}
\begin{tikzpicture}[
    x=1cm, y=1cm,
    >=Stealth,
    font=\sffamily\scriptsize,
    regnode/.style={
        draw=black!85, rectangle, rounded corners=3pt, minimum width=1.1cm, minimum height=0.75cm,
        line width=1.1pt, fill=white, font=\bfseries\scriptsize, align=center
    },
    block/.style={
        draw=black!85, rectangle, rounded corners=3pt, line width=1.1pt, fill=white,
        font=\bfseries\scriptsize, align=center
    },
    ctrlbox/.style={
        draw=red!80!black, rectangle, rounded corners=7pt, line width=1.4pt,
        fill=red!5, align=center
    },
    membox/.style={
        draw=cyan!75!black, rectangle, rounded corners=6pt, line width=1.3pt,
        fill=cyan!6, align=center
    },
    databus/.style={->, line width=1.3pt, draw=black!90},
    databus_line/.style={line width=1.3pt, draw=black!90},
    ctrlbus/.style={->, line width=1.0pt, draw=red!80!black, dashed},
    ctrlbus_line/.style={line width=1.0pt, draw=red!80!black, dashed},
    bidir/.style={<->, line width=1.2pt, draw=teal!80!black}
]

    \foreach \i/\x in {0/0.6, 1/2.0, 2/3.4, 3/4.8, 4/6.2, 5/7.6, 6/9.0, 7/10.4} {
        \node[regnode] (R\i) at (\x, 12.0) {\textbf{R\i}};
        \coordinate (R\i_in) at (\x - 0.25, 12.38);
        \coordinate (R\i_we) at (\x + 0.25, 12.38);
        \coordinate (R\i_out) at (\x, 11.62);
    }

    \draw[fill=gray!12, draw=black!85, line width=1.2pt]
        (-0.2, 10.6) -- (11.2, 10.6) -- (10, 9.7) -- (1.1, 9.7) -- cycle;
    \node[align=center] at (5.5, 10.15) {\textbf{\normalsize MUX 1}};

    \foreach \i in {0,...,7} {
        \draw[databus] (R\i_out) -- (R\i_out |- 0, 10.6);
    }

    \draw[fill=gray!12, draw=black!85, line width=1.1pt]
        (1.3, 6.9) -- (3.1, 6.9) -- (3.6, 5.9) -- (0.8, 5.9) -- cycle;
    \node[align=center] at (2.2, 6.4) {\textbf{\scriptsize DEST-DMUX}};

    \node[regnode, fill=blue!6, draw=blue!75!black, minimum width=1.2cm, minimum height=0.8cm]
        (ACCA) at (1.3, 4.9) {\textbf{ACCA}};
    \node[regnode, fill=blue!6, draw=blue!75!black, minimum width=1.2cm, minimum height=0.8cm]
        (ACCB) at (3.1, 4.9) {\textbf{ACCB}};

    \draw[databus] (1.3, 5.9) -- (ACCA.north);
    \draw[databus] (3.1, 5.9) -- (ACCB.north);
    \node[font=\tiny\bfseries, fill=white, inner sep=1pt] at (1.05, 5.55) {[0]};
    \node[font=\tiny\bfseries, fill=white, inner sep=1pt] at (3.35, 5.55) {[1]};

    \draw[fill=blue!8, draw=blue!85!black, line width=1.3pt]
        (0.5, 3.6) -- (1.8, 3.6) -- (2.2, 3.2) -- (2.6, 3.6) -- (3.9, 3.6) -- (3.2, 2.0) -- (1.2, 2.0) -- cycle;
    \node[align=center, text=blue!90!black] at (2.2, 2.7) {\textbf{\large ALU}};

    \draw[databus] (ACCA.south) -- (1.3, 3.6);
    \draw[databus] (ACCB.south) -- (3.1, 3.6);
    \node[font=\tiny\bfseries, left] at (1.3, 4.1) {A [15:0]};
    \node[font=\tiny\bfseries, right] at (3.1, 4.1) {B [15:0]};

    \draw[fill=yellow!18, draw=orange!80!black, line width=1.3pt]
        (5.6, 5.8) -- (7.2, 5.8) -- (7.8, 4.4) -- (5.0, 4.4) -- cycle;
    \node[align=center, text=orange!90!black] at (6.4, 5.25) {\textbf{\normalsize MUX 2}};

    \node[font=\tiny\bfseries, fill=white, draw=orange!80!black, rounded corners=1pt, inner sep=1.2pt] at (5.6, 4.68) {ALU};
    \node[font=\tiny\bfseries, fill=white, draw=orange!80!black, rounded corners=1pt, inner sep=1.2pt] at (6.4, 4.68) {REG};
    \node[font=\tiny\bfseries, fill=white, draw=orange!80!black, rounded corners=1pt, inner sep=1.2pt] at (7.2, 4.68) {MEM};

    \node[block, minimum width=3.0cm, minimum height=0.8cm, fill=teal!6, draw=teal!80!black]
        (FLAGS) at (6.4, 1.2) {\textbf{FREG [15:0]}\\[1pt]\scriptsize ZF, CF, OF, EM, \ldots};

    \node[ctrlbox, minimum width=5.4cm, minimum height=5.0cm]
        (CTRL) at (14.1, 6.7) {};

    \node[font=\bfseries\normalsize, text=red!90!black] at (14.1, 8.7) {CONTROL UNIT};

    \node[draw=red!40!black, fill=white, rounded corners=2pt, font=\tiny\bfseries, inner sep=2pt, right]
        at (11.5, 8.3) {IR [15:0]};
    \node[draw=red!70!black, fill=white, rounded corners=2pt, font=\tiny\bfseries, inner sep=2pt, text=red!80!black, right]
        at (11.5, 7.3) {CTRL ALU/ACC};
    \node[draw=red!70!black, fill=white, rounded corners=2pt, font=\tiny\bfseries, inner sep=2pt, text=red!80!black, right]
        at (11.5, 6.2) {SEL\_WB [1:0]};
    \node[draw=teal!70!black, fill=white, rounded corners=2pt, font=\tiny\bfseries, inner sep=2pt, text=teal!90!black, right]
        at (11.5, 5.2) {FLAGS};
    \node[draw=orange!80!black, fill=white, rounded corners=2pt, font=\tiny\bfseries, inner sep=2pt, text=orange!90!black, right]
        at (11.5, 4.5) {MEM/IMM};

    \node[membox, minimum width=5.4cm, minimum height=2.6cm]
        (MEM) at (14.1, 2.1) {\textbf{\normalsize MEMORY}\\[3pt]\tiny RAM / ROM / DEVICES / ELSE};

    \draw[databus, draw=blue!80!black] (12.6, 4.2) -- (12.6, 3.4) node[midway, fill=white, inner sep=1pt, font=\tiny\bfseries, text=blue!80!black] {ADDR};
    \draw[databus] (14.1, 4.2) -- (14.1, 3.4) node[midway, fill=white, inner sep=1pt, font=\tiny\bfseries] {WR\_DATA};
    \draw[bidir, draw=red!80!black] (15.6, 3.4) -- (15.6, 4.2) node[midway, fill=white, inner sep=1pt, font=\tiny\bfseries, text=red!80!black] {RD / WR};

    \draw[databus_line] (5.5, 9.7) -- (5.5, 8.3);
    \filldraw (5.5, 8.3) circle (2pt);

    \draw[databus] (5.5, 8.3) -| (2.2, 6.9) node[pos=0.25, above, font=\tiny\bfseries] {[15:0]};

    \draw[databus] (5.5, 8.3) -- (11.4, 8.3) node[pos=0.45, above, font=\tiny\bfseries] {REG\_DATA [15:0]};

    \filldraw (4.6, 8.3) circle (1.8pt);
    \draw[databus_line] (4.6, 8.3) -- (4.6, 3.6);
    \draw[databus_line] (4.6, 3.6) -- (5.4, 3.6);
    \draw[databus_line] (5.4, 3.6) arc[start angle=180, end angle=0, radius=0.2] -- (6.4, 3.6);
    \draw[databus] (6.4, 3.6) -- (6.4, 4.4);
    \node[font=\tiny\bfseries, fill=white, inner sep=1.2pt, below] at (5.0, 3.55) {[15:0]};

    \draw[databus] (3.4, 2.5) -- (5.6, 2.5) -- (5.6, 4.4);
    \node[font=\tiny\bfseries, fill=white, inner sep=1pt] at (4.5, 2.75) {ALU\_RES [15:0]};

    \draw[databus_line] (11.4, 4.5) -- (10.3, 4.5) -- (10.3, 3.6) -- (9.8, 3.6);
    \draw[databus_line] (9.8, 3.6) arc[start angle=0, end angle=180, radius=0.2] -- (7.2, 3.6);
    \draw[databus] (7.2, 3.6) -- (7.2, 4.4);
    \node[font=\tiny\bfseries, fill=white, inner sep=1pt] at (10.8, 4.7) {MEM/IMM [15:0]};

    \draw[->, line width=1.1pt, draw=teal!80!black] (2.2, 2.0) |- (FLAGS.west);
    \node[font=\tiny\bfseries, fill=white, inner sep=1pt, text=teal!90!black] at (3.5, 1.2) {FLAGS};

    \draw[bidir] (7.9, 1.2) -- (9.6, 1.2) -- (9.6, 5.2) -- (11.4, 5.2);
    \node[font=\tiny\bfseries, fill=white, inner sep=1pt, text=teal!90!black] at (10.4, 5.45) {FLAGS};

    \draw[databus_line] (6.4, 5.8) -- (6.4, 9.0) -- (5.7, 9.0)
                        arc[start angle=0, end angle=180, radius=0.2]
                        -- (-0.8, 9.0) -- (-0.8, 13.5) -- (10.9, 13.5);
    \node[font=\tiny\bfseries, fill=black!85, text=white, inner sep=2.5pt, rounded corners=2pt]
        at (2.2, 9.0) {WB\_DATA [15:0]};

    \foreach \i in {0,...,7} {
        \draw[databus] (R\i_in |- 0, 13.5) -- (R\i_in);
        \filldraw (R\i_in |- 0, 13.5) circle (1.8pt);
    }

    \draw[ctrlbus] (12.8, 9.2) |- (10.6 , 10.15);
    \node[font=\tiny\bfseries, text=red!80!black, fill=white, inner sep=1pt] at (12.8, 9.7) {SEL\_REG [2:0]};

    \draw[ctrlbus] (11.4, 6.2) -| (8.5, 5.3) -- (7.4, 5.3);
    \node[font=\tiny\bfseries, text=red!80!black, fill=white, inner sep=1pt] at (9.8, 6.2) {SEL\_WB [1:0]};

    \draw[ctrlbus_line] (11.4, 7.3) -- (6.6, 7.3)
                        arc[start angle=0, end angle=180, radius=0.2]
                        -- (4.8, 7.3)
                        arc[start angle=0, end angle=180, radius=0.2]
                        -- (4.1, 7.3);

    \draw[ctrlbus_line] (4.1, 7.3) -- (4.1, 4.9);
    \filldraw[red!80!black] (4.1, 7.3) circle (1.4pt);

    \draw[ctrlbus] (4.1, 6.4) -- (3.35, 6.4);
    \filldraw[red!80!black] (4.1, 6.4) circle (1.4pt);
    \node[font=\tiny\bfseries, text=red!80!black, fill=white, inner sep=1pt, above] at (4.1, 6.45) {SEL\_ACC};

    \draw[ctrlbus] (4.1, 4.9) -- (ACCB.east);
    \node[font=\tiny\bfseries, text=red!80!black, fill=white, inner sep=1pt, above] at (4.15, 5.1) {LD\_B};

    \draw[ctrlbus_line] (4.1, 7.3) -- (2.4, 7.3)
                        arc[start angle=0, end angle=180, radius=0.2]
                        -- (0.3, 7.3);

    \draw[ctrlbus] (0.2, 7.3) -- (0.2, 4.9) -- (ACCA.west);
    \node[font=\tiny\bfseries, text=red!80!black, fill=white, inner sep=1pt, left] at (0.4, 5.2) {LD\_A};

    \draw[ctrlbus] (0.2, 4.9) -- (0.2, 2.8) -- (0.8, 2.8);
    \node[font=\tiny\bfseries, text=red!80!black, fill=white, inner sep=1pt, left] at (0.5, 3.3) {ALU\_OP [3:0]};

    \draw[ctrlbus_line] (15.2, 9.2) -- (15.2, 12.8) -- (-0.2, 12.8);
    \node[font=\tiny\bfseries, text=red!80!black, fill=white, inner sep=1pt] at (13.0, 12.8) {REG\_WE [7:0]};
    \foreach \i in {0,...,7} {
        \draw[ctrlbus] (R\i_we |- 0, 12.8) -- (R\i_we);
        \filldraw[red!80!black] (R\i_we |- 0, 12.8) circle (1.4pt);
    }

\end{tikzpicture}
\end{center}
\end{landscape}
\clearpage

\section{Программно видимые регистры и состояние процессора}
Архитектурное состояние процессора LIM M2 представлено набором ортогональных регистров, сгруппированных по функциональному назначению:

\begin{longtable}{>{\ttfamily\bfseries}p{25mm}p{28mm}p{93mm}}
\toprule
Регистр & Разрядность & Архитектурная роль и особенности \\
\midrule\endhead
R0--R7 & 16 бит & Регистры общего назначения (GPR). Селектор MUX 1 обеспечивает сверхбыстрое бесконфликтное считывание; фиксация результата происходит по стробу с шины WB\_DATA. \\
SREG & 24 бита & Системный регистр стека (Stack Register). Поддерживает эффективное аппаратное отслеживание глубины стековых кадров с прямой 24-битной адресацией стекового пространства (до 16~Мбайт). \\
PTREG & 24 бита & Высокоточный регистр-указатель (Pointer Register). Обеспечивает прямое линейное обращение к памяти объёмом до 16~Мбайт с аппаратным автоинкрементом по флагу AIP. \\
FREG & 16 бит & Регистр состояния и управления (Processor Status \& Control Register). Объединяет флаги арифметических результатов и биты конфигурации режимов ядра. \\
\bottomrule
\end{longtable}

\subsection{Топология регистра флагов и управления (FREG)}
Регистр \reg{FREG} является ключевым элементом контроля математического тракта и системных режимов:

\begin{tabularx}{\textwidth}{>{\bfseries}p{14mm}>{\ttfamily}p{23mm}X}
\toprule
Бит & Мнемоника & Функциональное описание \\
\midrule
0    & ZF   & Zero Flag: истинен при нулевом результате операции. \\
1    & CF   & Carry Flag: аппаратный флаг переноса / заёма в арифметике. \\
2    & OF   & Overflow Flag: индикатор знакового переполнения дополнительного кода. \\
3    & EM   & Error Math: аппаратная ловушка математических исключений (деление на ноль). \\
4--5 & CMP  & Compare Result: результат аппаратного компаратора (00 --- сброс/нет сравнения/ошибка, 01 --- меньше, 10 --- равно, 11 --- больше). \\
6--7 & CLK  & Clock Mode: выбор профиля делителя тактовой частоты ядра. \\
8    & PIE  & Program Interrupt Enable: глобальное разрешение программных прерываний. \\
9    & MIE  & Master Interrupt Enable: разрешение маскируемых внешних прерываний. \\
10   & DWR  & Disable Wait for Ready: форсированный безакцептный обмен без ожидания шины. \\
11   & M24  & 24-bit Mode Enable: активация расширенного 24-битного адресного режима. \\
12   & AIP  & Auto-Increment PTREG: разрешение аппаратного инкремента указателя \reg{PTREG} при обращении. \\
13   & DBIT & Displaced Bit: буфер вытесненного бита сдвигового узла. \\
14   & SMOD & Shift Mode: управление режимом сдвига (0 --- заполнение нулями, 1 --- циклический проход через DBIT). \\
15   & INTR & In Interrupt: флаг нахождения процессора в прерывании (при 1 новые прерывания аппаратно заблокированы). \\
\bottomrule
\end{tabularx}

\section{Формат машинных инструкций}
Инструкции процессора LIM M2 обладают компактным, легко декодируемым форматом. 
Базовый опкод занимает старшие 5 бит, обеспечивая мгновенную диспетчеризацию в микропрограммном ПЗУ управления:

\begin{center}
\begin{tikzpicture}[x=0.9cm,y=0.72cm]
  \draw[thick,fill=limblue] (0,0) rectangle (5,1.25);
  \draw[thick,fill=limcyan] (5,0) rectangle (8,1.25);
  \draw[thick,fill=limlight] (8,0) rectangle (17,1.25);
  \node[text=white,font=\bfseries,align=center] at (2.5,0.62) {OPCODE\\5 бит};
  \node[text=white,font=\bfseries,align=center] at (6.5,0.62) {CFG\\3 бита};
  \node[font=\bfseries,align=center] at (12.5,0.62) {Операнды и аргументы\\динамическая длина};
  \node[anchor=south west,font=\small] at (0,1.32) {$W-1$};
  \node[anchor=south,font=\small] at (5,1.32) {$W-6$};
  \node[anchor=south,font=\small] at (8,1.32) {$W-9$};
  \node[anchor=south east,font=\small] at (17,1.32) {0};
\end{tikzpicture}
\end{center}

Поле \field{cfg} (3 бита) определяет специализированную конфигурацию (до 8 ортогональных режимов операции), что позволяет одной базовой команде покрывать широкий спектр адресаций без раздувания пространства опкодов.

\section{Сводная таблица команд ядра LIM M2}
\small
\begin{longtable}{r>{\ttfamily}p{13mm}>{\ttfamily}p{19mm}p{28mm}p{75mm}}
\toprule
№ & Код & Мнемоника & Класс & Назначение операции \\
\midrule\endhead
0  & 00000 & NOP  & Управление & Пустая операция и сброс остаточного регистра АЛУ \\
1  & 00001 & ADD  & Арифметика & Сложение с автокаскадированием переноса через остаток АЛУ \\
2  & 00010 & SUB  & Арифметика & Вычитание с автокаскадированием заёма через остаток АЛУ \\
3  & 00011 & MUL  & Арифметика & Умножение с фиксацией старших разрядов в остатке АЛУ \\
4  & 00100 & DIV  & Арифметика & Деление с фиксацией остатка во внутреннем регистре АЛУ \\
5  & 00101 & LMR  & Память     & Загрузка из памяти/ввода в регистр (конфиг. источника 0--3, размер 8/16 бит) \\
6  & 00110 & LRM  & Память     & Запись данных из регистра в память/вывод (конфиг. назначения 0--3, размер 8/16 бит) \\
7  & 00111 & LRR  & Пересылка  & Прямое межрегистровое копирование через обходной MUX \\
8  & 01000 & CMP  & Сравнение  & Аппаратное сравнение операндов без разрушения данных \\
9  & 01001 & JMP  & Переход    & Безусловный переход (прямой, косвенный, 24-битный) \\
10 & 01010 & JCR  & Переход    & Условный переход по состоянию флага переноса CF \\
11 & 01011 & JNZ  & Переход    & Условный переход при ненулевом результате (ZF = 0) \\
12 & 01100 & JZ   & Переход    & Условный переход при нулевом результате (ZF = 1) \\
13 & 01101 & AND  & Логика     & Побитовое логическое умножение (И) \\
14 & 01110 & OR   & Логика     & Побитовое логическое сложение (ИЛИ) \\
15 & 01111 & NAND & Логика     & Побитовая инверсия логического И \\
16 & 10000 & NOR  & Логика     & Побитовая инверсия логического ИЛИ \\
17 & 10001 & NOT  & Логика     & Побитовое инвертирование операнда (дополнение до 1) \\
18 & 10010 & XOR  & Логика     & Побитовое исключающее ИЛИ \\
19 & 10011 & XNOR & Логика     & Побитовая эквивалентность \\
20 & 10100 & INC  & Арифметика & Аппаратный инкремент операнда на единицу \\
21 & 10101 & DEC  & Арифметика & Аппаратный декремент операнда на единицу \\
22 & 10110 & PUSH & Стек       & Размещение регистра в стековом пространстве \\
23 & 10111 & POP  & Стек       & Извлечение вершины стека в целевой регистр \\
24 & 11000 & BSL  & Сдвиги     & Баррельный сдвиг влево с настраиваемым вводом бит \\
25 & 11001 & BSR  & Сдвиги     & Баррельный сдвиг вправо с контролем флага DBIT \\
26 & 11010 & MSP  & Система    & Манипуляция специальными регистрами (FREG, PTREG) \\
27 & 11011 & RSV27 & Резерв    & Зарезервировано для векторных/SIMD инструкций \\
28 & 11100 & RSV28 & Резерв    & Зарезервировано для расширенной адресации \\
29 & 11101 & RSV29 & Резерв    & Зарезервировано для операций с плавающей точкой \\
30 & 11110 & RSV30 & Резерв    & Зарезервировано для аппаратной криптографии \\
31 & 11111 & RSV31 & Резерв    & Зарезервировано для расширения системы команд \\
\bottomrule
\end{longtable}
\normalsize

\section{Детальная спецификация системы команд}
Ниже представлены исчерпывающие паспорта всех 32 базовых опкодов процессора LIM M2.

\InstructionPage{NOP}{00000}{Управление}{Аппаратная пустая операция и сброс арифметической цепочки}{%
\small
Сигнатура: NOP \\
Операция:  PC <- PC + 1, ALU\_RESIDUAL <- 0, FREG.CF <- 0, FREG.OF <- 0 \\[2pt]
Пример:    NOP              ; Сброс цепочки переноса АЛУ
}{Операнды отсутствуют; поле \field{cfg} зарезервировано для калибровки задержек.}{Сбрасывает флаги переноса (CF), переполнения (OF) и заёма каскадной цепочки АЛУ; остальные флаги сохраняются.}{Служит для тактовой синхронизации конвейера и аппаратного сброса (Flush) внутреннего остаточного регистра АЛУ, разрывая цепочку автопереноса перед новыми независимыми вычислениями.}
\InstructionPage{ADD}{00001}{Арифметика}{Сложение с автоматическим каскадированием переноса}{%
\small
Сигнатура: ADD arg0(REG), arg1(REG), arg2(REG) : arg0 = arg1 + arg2 \\
Операция:  sum <- arg1 + arg2 + ALU\_RESIDUAL \\
           arg0 <- sum[15:0], ALU\_RESIDUAL <- (sum > 0xFFFF ? 1 : 0) \\[2pt]
Пример:    ADD R0, R1, R2   ; R0 = R1 + R2 (регистры R0--R7)
}{Регистры R0--R7, аккумуляторы ACCA/ACCB или непосредственные операнды.}{Взводит флаг переноса CF и переполнения OF при превышении $2^{16}-1$; обновляет ZF по значению dst.}{Поддерживает сквозную многобайтовую арифметику (32, 48, 64+ бит). Если результат суммы не помещается в 16 бит ($>2^{16}-1$), перенос (+1) фиксируется во внутреннем регистре АЛУ и автоматически суммируется при следующей операции ADD без отдельных инструкций ADC. Младшие 16 бит всегда детерминированно сохраняются в dst. Цепочка сбрасывается инструкцией NOP либо перезаписывается новой математической операцией.}
\InstructionPage{SUB}{00010}{Арифметика}{Вычитание с автоматическим каскадированием заёма}{%
\small
Сигнатура: SUB arg0(REG), arg1(REG), arg2(REG) : arg0 = arg1 - arg2 \\
Операция:  diff <- arg1 - arg2 - ALU\_RESIDUAL \\
           arg0 <- diff[15:0], ALU\_RESIDUAL <- (diff < 0 ? 1 : 0) \\[2pt]
Пример:    SUB R0, R1, R2   ; R0 = R1 - R2 (регистры R0--R7)
}{Регистры R0--R7, константы или аккумуляторы ACCA/ACCB.}{При разности $< 0$ взводит флаг заёма (Borrow / CF) и OF; обновляет ZF.}{Реализует непрерывное многословное вычитание. При возникновении отрицательной разности ($< 0$) заём ($-1$) запоминается во внутреннем регистре АЛУ, и в следующей операции SUB АЛУ автоматически вычитает 1 из старшего слова. Результат младшего слова детерминированно фиксируется в dst. Для начала новой независимой разности цепочка прерывается инструкцией NOP.}
\InstructionPage{MUL}{00011}{Арифметика}{Целочисленное умножение с сохранением старшей части в регистре АЛУ}{%
\small
Сигнатура: MUL arg0(REG), arg1(REG), arg2(REG) : arg0 = arg1 * arg2 \\
Операция:  prod <- arg1 * arg2 \\
           arg0 <- prod[15:0], ALU\_RESIDUAL <- prod[31:16] \\[2pt]
Пример:    MUL R0, R1, R2   ; R0 = R1 * R2 (старшие 16 бит в ALU\_RESIDUAL)
}{Регистры R0--R7, аккумуляторы ACCA/ACCB.}{Обновляет ZF; взводит OF/CF, если старшие разряды произведения (prod[31:16]) не равны 0.}{При умножении младшие 16 бит произведения направляются в целевой регистр dst, а старшие 16 бит сохраняются во внутреннем служебном регистре АЛУ (ALU Residual). При следующей операции сложения сохранённый остаток автоматически суммируется со следующим результатом (аппаратная поддержка MAC / многословных чисел). Старшую часть произведения можно также извлечь в регистр, выполнив сложение регистра с самим собой после очистки. Инструкция NOP очищает остаточный регистр АЛУ.}
\InstructionPage{DIV}{00100}{Арифметика}{Целочисленное деление с сохранением остатка во внутреннем регистре АЛУ}{%
\small
Сигнатура: DIV arg0(REG), arg1(REG), arg2(REG) : arg0 = arg1 / arg2 \\
Операция:  arg0 <- arg1 / arg2 \\
           ALU\_RESIDUAL <- arg1 \% arg2 \\[2pt]
Пример:    DIV R0, R1, R2   ; R0 = R1 / R2 (остаток modulo в ALU\_RESIDUAL)
}{Регистры R0--R7, аккумуляторы ACCA/ACCB.}{При делении на ноль взводится аппаратный флаг ошибки EM (Error Math) без зависания ядра; обновляет ZF.}{Результат целочисленного деления (частное) записывается в целевой регистр dst, а остаток от деления (modulo) автоматически сохраняется во внутреннем служебном регистре АЛУ. Сохранённый остаток может быть задействован в последующих операциях АЛУ либо сброшен командой NOP.}
\InstructionPage{LMR}{00101}{Память}{Загрузка данных из памяти и внешних устройств в регистры}{%
\small
Сигнатура: LMR arg0(REG), <SRC\_MODE> : arg0 = MEM[...] \\
Режимы:    SRC=0: LMR R0, [PC+]          ; Imm Data (1--2 байта) \\
           SRC=1: LMR R0, [[PC+]]        ; Imm Pointer \\
           SRC=2: LMR R0, [PTREG]        ; Системный указатель (AIP) \\
           SRC=3: LMR R0(R\_hi), R1(R\_lo) ; 24-бит адрес из пары \\
Операция:  arg0 <- MEM[EA] (старшие биты адреса перезаписываются) \\[2pt]
Пример:    LMR R0, R1       ; EA=(R0<<16)|R1, чтение MEM[EA] в R0
}{%
\textbf{Октет 0:} бит $[8]$ (\field{SZ}) задаёт разрядность ($0 = 1$ байт, $1 = 2$ байта / слово); биты $[10:9]$ (\field{SRC}) задают источник данных ($0$ --- данные сразу за инструкцией в памяти, $1$ --- непосредственный указатель за инструкцией, $2$ --- косвенно по 24-битному указателю \reg{PTREG}, $3$ --- косвенно по паре пользовательских регистров).\\
\textbf{Октет 1:} биты $[7:5]$ задают регистр \reg{Arg1} (\reg{R\_hi} / приёмник), биты $[4:2]$ задают регистр \reg{Arg2} (\reg{R\_lo}). В режиме $\field{SRC} = 3$ два 16-битных регистра формируют 32-битное число, из которого для адресации используются 24 бита: \reg{Arg1} хранит старшие биты адреса, а \reg{Arg2} --- младшие 16 бит. При этом \reg{Arg1} одновременно является целевым регистром: после выборки из памяти старшие биты адреса в нём перезаписываются считанными данными.%
}{Все флаги регистра состояния \reg{FREG} сохраняются без изменений.}{%
Инструкция \texttt{LMR} является фундаментальным шлюзом ввода внешней информации в регистровый файл процессора LIM M2 ($R0$--$R7$). Архитектура поддерживает четыре ортогональных режима выборки и полную 24-битную адресацию:\\
1. В режиме $\field{SRC} = 0$ процессор считывает константу непосредственно из потока кода ($1$ или $2$ байта) со смещением счетчика команд \reg{PC}.\\
2. В режиме $\field{SRC} = 1$ за инструкцией считывается абсолютный указатель, по которому загружаются данные.\\
3. В режиме $\field{SRC} = 2$ чтение производится по 24-битному физическому адресу из системного регистра \reg{PTREG}. При активном флаге автоинкремента \texttt{FREG.AIP} (бит 12) регистр \reg{PTREG} автоматически инкрементируется на размер выборки ($+1$ или $+2$), обеспечивая аппаратный потоковый ввод массивов данных за один такт без накладных расходов.\\
4. В режиме $\field{SRC} = 3$ пользователь выбирает два 16-битных регистра (\reg{Arg1} и \reg{Arg2}), которые вместе образуют 32-битное число, из которого на практике адресации используются младшие 24 бита. Пользователь указывает, что старшие биты адреса находятся в \reg{Arg1}, а младшие 16 бит адреса --- в \reg{Arg2}. В этом режиме регистр \reg{Arg1} со старшими битами адреса также становится регистром назначения, куда сохраняются считанные данные (то есть старшие биты адреса в \reg{Arg1} перезаписываются данными, в то время как \reg{Arg2} сохраняет младшие биты адреса).%
}
\InstructionPage{LRM}{00110}{Память}{Запись данных из регистров в память и внешние устройства}{%
\small
Сигнатура: LRM arg0(REG), <DST\_MODE> : MEM[...] = arg0 \\
Режимы:    DST=0: LRM R0, [PC+]          ; Inline память за кодом \\
           DST=1: LRM R0, [[PC+]]        ; Imm Pointer \\
           DST=2: LRM R0, [PTREG]        ; Системный указатель (AIP) \\
           DST=3: LRM R0(R\_hi), R1(R\_lo) ; 24-бит адрес из пары \\
Операция:  MEM[EA] <- arg0[15:0] (или младший байт при SZ=0) \\[2pt]
Пример:    LRM R0, R1       ; EA=(R0<<16)|R1, запись R0 в MEM[EA]
}{%
\textbf{Октет 0:} бит $[8]$ (\field{SZ}) задаёт разрядность ($0 = 1$ байт, $1 = 2$ байта / слово); биты $[10:9]$ (\field{DST}) задают целевой приёмник данных ($0$ --- запись сразу за кодом инструкции, $1$ --- запись по указателю за инструкцией, $2$ --- косвенно по 24-битному указателю \reg{PTREG}, $3$ --- косвенно по паре регистров пользователя).\\
\textbf{Октет 1:} биты $[7:5]$ задают регистр-источник \reg{Arg1} (\reg{Rs} / \reg{R\_hi}), биты $[4:2]$ задают регистр \reg{Arg2} (\reg{R\_lo}). В режиме $\field{DST} = 3$ два регистра формируют 32-битное число, из которого используются 24 бита адреса (\reg{Arg1} --- старшие биты адреса и источник данных, \reg{Arg2} --- младшие 16 бит). Уничтожение существующих данных по адресу записи является ответственностью программиста.%
}{Все флаги регистра состояния \reg{FREG} сохраняются без изменений.}{%
Инструкция \texttt{LRM} является фундаментальным механизмом выгрузки регистровых данных процессора LIM M2 ($R0$--$R7$) в оперативную память и внешние порты периферии. Архитектура поддерживает четыре ортогональных режима адресации:\\
1. В режиме $\field{DST} = 0$ процессор записывает $1$ или $2$ байта непосредственно в память вслед за инструкцией (\reg{PC}) со сдвигом счетчика команд. Режим применим для динамических буферов и самомодификации кода; перезапись данных лежит на ответственности программиста.\\
2. В режиме $\field{DST} = 1$ за инструкцией считывается 16-битный абсолютный адрес, по которому сохраняются данные.\\
3. В режиме $\field{DST} = 2$ запись производится по 24-битному системному адресу регистра \reg{PTREG}. При активном флаге \texttt{FREG.AIP} (бит 12) регистр \reg{PTREG} автоматически инкрементируется на $+1$ или $+2$, обеспечивая аппаратный потоковый DMA-вывод массивов в память без программных циклов.\\
4. В режиме $\field{DST} = 3$ используется пара 16-битных регистров: \reg{Arg1} и \reg{Arg2} образуют 32-битное составное число, из которого младшие 24 бита задают физический адрес ($\text{Addr}[23:0]$). Данные из \reg{Arg1} записываются по вычисленному адресу, обеспечивая прямой доступ ко всему адресному пространству 16 МБ.%
}
\InstructionPage{LRR}{00111}{Пересылка}{Межрегистровое перемещение (копирование из точки Б в точку А)}{%
\small
Сигнатура: LRR arg0(REG), arg1(REG) : arg0 = arg1 \\
Операция:  arg0 <- arg1 (прямое копирование из точки Б в точку А) \\[2pt]
Пример:    LRR R0, R1       ; R0 = R1 (значение R1 копируется в R0)
}{%
\textbf{Октет 0:} код операции $00111_2$ (\texttt{0x07}), биты $[10:8]$ зарезервированы ($000_2$).\\
\textbf{Октет 1:} биты $[7:5]$ задают целевой регистр-приёмник \reg{Arg1} (\reg{Rd}, точка А, куда сохраняются данные); биты $[4:2]$ задают регистр-источник \reg{Arg2} (\reg{Rs}, точка Б, откуда берутся данные); биты $[1:0]$ содержат $00_2$. При исполнении значение \reg{Rs} передаётся в \reg{Rd} за 1 микротакт.%
}{Все флаги регистра состояния \reg{FREG} сохраняются без изменений.}{%
Инструкция \texttt{LRR} является фундаментальной операцией пересылки данных между регистрами общего назначения процессора ($R0$--$R7$). Она реализует передачу из источника (точка Б) в приёмник (точка А). Пересылка выполняется через прямой аппаратный обходной тракт \textbf{MUX 1} $\to$ \textbf{MUX 2} минуя АЛУ, что исключает блокировку вычислительного тракта и не влияет на флаги регистра \reg{FREG}.%
}
\InstructionPage{CMP}{01000}{Сравнение}{Аппаратное сравнение операндов}{%
\small
Сигнатура: CMP arg0(REG), arg1(REG) : compare(arg0, arg1) \\
Операция:  if arg0 < arg1 then FREG.CMP <- 01b (меньше) \\
           if arg0 == arg1 then FREG.CMP <- 10b (равно) \\
           if arg0 > arg1 then FREG.CMP <- 11b (больше) \\
           (состояние 00b --- сброс / не сравнивалось / ошибка) \\[2pt]
Пример:    CMP R0, R1       ; 01: R0<R1, 10: R0==R1, 11: R0>R1 (00: сброс/ошибка)
}{Сравнение регистров без изменения данных.}{Обновляет двухбитное поле FREG.CMP (биты [5:4]), а также флаги ZF, CF, OF.}{Позволяет выполнять быстрые условные переходы без искажения регистров.}
\InstructionPage{JMP}{01001}{Переход}{Безусловный переход по адресу}{%
\small
Сигнатура: JMP target(IMM|REG) : PC = target \\
Операция:  PC <- target \\[2pt]
Пример:    JMP R0           ; Безусловный переход по адресу в R0
}{Поддерживает 16-битный внутристраничный или полный 24-битный переход через \reg{PTREG}.}{Флаги сохраняются без изменений.}{Атомарная загрузка Program Counter за минимальное число микротактов.}
\InstructionPage{JCR}{01010}{Переход}{Условный переход по флагу переноса}{%
\small
Сигнатура: JCR target(IMM|REG) : if FREG.CF == 1 then PC = target \\
Операция:  if FREG.CF == 1 then PC <- target \\[2pt]
Пример:    JCR R2           ; Переход по адресу R2 при наличии переноса
}{Целевой адрес задаётся смещением либо абсолютным значением.}{Флаги читаются, но не изменяются.}{Ключевой примитив для построения длинной многословной арифметики.}
\InstructionPage{JNZ}{01011}{Переход}{Условный переход при неравенстве нулю}{%
\small
Сигнатура: JNZ target(IMM|REG) : if FREG.ZF == 0 then PC = target \\
Операция:  if FREG.ZF == 0 then PC <- target \\[2pt]
Пример:    JNZ R3           ; Переход по адресу R3 при ZF == 0 (не ноль)
}{Высокоэффективный переход для организации циклов и счетчиков.}{Флаги не изменяются.}{Минимизированное время принятия решения в микрокоде.}
\InstructionPage{JZ}{01100}{Переход}{Условный переход при равенстве нулю}{%
\small
Сигнатура: JZ target(IMM|REG) : if FREG.ZF == 1 then PC = target \\
Операция:  if FREG.ZF == 1 then PC <- target \\[2pt]
Пример:    JZ R4            ; Переход по адресу R4 при ZF == 1 (равно 0)
}{Отрабатывает условие успешного завершения или совпадения.}{Флаги не изменяются.}{Оптимизирован для алгоритмов поиска и сравнения строк.}
\InstructionPage{AND}{01101}{Логика}{Побитовое логическое И}{%
\small
Сигнатура: AND arg0(REG), arg1(REG) : arg0 = arg0 \& arg1 \\
Операция:  arg0 <- arg0 AND arg1 \\[2pt]
Пример:    AND R0, R1   ; R0 = R0 \& R1
}{Регистровые операнды (2 аргумента).}{Обновляет флаг нуля ZF; сбрасывает флаги переполнения.}{Идеально для наложения битовых масок.}
\InstructionPage{OR}{01110}{Логика}{Побитовое логическое ИЛИ}{%
\small
Сигнатура: OR arg0(REG), arg1(REG) : arg0 = arg0 | arg1 \\
Операция:  arg0 <- arg0 OR arg1 \\[2pt]
Пример:    OR R0, R1    ; R0 = R0 | R1
}{Регистровые операнды (2 аргумента).}{Обновляет флаг ZF; флаги переноса детерминированы.}{Быстрое объединение управляющих битов.}
\InstructionPage{NAND}{01111}{Логика}{Побитовое логическое И-НЕ}{%
\small
Сигнатура: NAND arg0(REG), arg1(REG) : arg0 = \textasciitilde(arg0 \& arg1) \\
Операция:  arg0 <- NOT (arg0 AND arg1) \\[2pt]
Пример:    NAND R0, R1  ; R0 = \textasciitilde(R0 \& R1)
}{Универсальный логический базис (2 аргумента).}{Обновляет флаг ZF.}{Позволяет сократить код при синтезе нестандартных логических функций.}
\InstructionPage{NOR}{10000}{Логика}{Побитовое логическое ИЛИ-НЕ}{%
\small
Сигнатура: NOR arg0(REG), arg1(REG) : arg0 = \textasciitilde(arg0 | arg1) \\
Операция:  arg0 <- NOT (arg0 OR arg1) \\[2pt]
Пример:    NOR R0, R1   ; R0 = \textasciitilde(R0 | R1)
}{Полноценная инверсная логика (2 аргумента).}{Обновляет флаг ZF.}{Расширяет набор быстрых логических трансформаций.}
\InstructionPage{NOT}{10001}{Логика}{Побитовое инвертирование (дополнение до 1)}{%
\small
Сигнатура: NOT arg0(REG) : arg0 = \textasciitilde arg0 \\
Операция:  arg0 <- NOT arg0 \\[2pt]
Пример:    NOT R0       ; R0 = \textasciitilde R0 (1-байтовая инструкция)
}{Однобайтовая инструкция ([7:3] опкод 10001, [2:0] Rd). Второй байт не запрашивается из памяти.}{Обновляет флаг ZF.}{Выполняется за минимальное количество микротактов без чтения второго байта.}
\InstructionPage{XOR}{10010}{Логика}{Побитое исключающее ИЛИ}{%
\small
Сигнатура: XOR arg0(REG), arg1(REG) : arg0 = arg0 \textasciicircum\ arg1 \\
Операция:  arg0 <- arg0 XOR arg1 \\[2pt]
Пример:    XOR R0, R1   ; R0 = R0 \textasciicircum\ R1 (XOR R0, R0 обнуляет R0)
}{Регистровые операнды (2 аргумента).}{Обновляет флаг ZF.}{Базовая операция для алгоритмов контрольных сумм и шифрования.}
\InstructionPage{XNOR}{10011}{Логика}{Побитовая эквивалентность}{%
\small
Сигнатура: XNOR arg0(REG), arg1(REG) : arg0 = \textasciitilde(arg0 \textasciicircum\ arg1) \\
Операция:  arg0 <- NOT (arg0 XOR arg1) \\[2pt]
Пример:    XNOR R0, R1  ; R0 = \textasciitilde(R0 \textasciicircum\ R1)
}{Побитовое сравнение на идентичность (2 аргумента).}{Обновляет флаг ZF.}{Эффективная проверка соответствия битовых масок.}
\InstructionPage{INC}{10100}{Арифметика}{Инкремент операнда на единицу}{%
\small
Сигнатура: INC arg0(REG) : arg0 = arg0 + 1 \\
Операция:  arg0 <- arg0 + 1 \\[2pt]
Пример:    INC R0           ; R0 = R0 + 1 (инкремент регистра)
}{Быстрый шаг итератора цикла.}{Обновляет ZF, OF; флаг переноса CF сохраняется или обновляется по \field{cfg}.}{Оптимизирует обработку массивов.}
\InstructionPage{DEC}{10101}{Арифметика}{Декремент операнда на единицу}{%
\small
Сигнатура: DEC arg0(REG) : arg0 = arg0 - 1 \\
Операция:  arg0 <- arg0 - 1 \\[2pt]
Пример:    DEC R0           ; R0 = R0 - 1 (декремент регистра)
}{Шаг обратного счетчика.}{Обновляет ZF, OF.}{Компактная реализация декрементных циклов.}
\InstructionPage{PUSH}{10110}{Стек}{Сохранение слова в стек}{%
\small
Сигнатура: PUSH arg0(REG) : SREG -= 2, MEM[SREG] = arg0 \\
Операция:  SREG <- SREG - 2, MEM[SREG] <- arg0 \\[2pt]
Пример:    PUSH R0          ; Сохранить слово R0 в вершину стека
}{Операнд: любой из регистров R0--R7 или специальный регистр.}{Флаги не изменяются.}{Аппаратная поддержка вызова процедур и сохранения контекста.}
\InstructionPage{POP}{10111}{Стек}{Извлечение слова из вершины стека}{%
\small
Сигнатура: POP arg0(REG) : arg0 = MEM[SREG], SREG += 2 \\
Операция:  arg0 <- MEM[SREG], SREG <- SREG + 2 \\[2pt]
Пример:    POP R0           ; Извлечь слово из вершины стека в R0
}{Целевой регистр: R0--R7 или регистры состояния.}{Флаги не изменяются.}{Симметричное восстановление контекста прерываний и функций.}
\InstructionPage{BSL}{11000}{Сдвиги}{Баррельный сдвиг влево}{%
\small
Сигнатура: BSL arg0(REG), arg1(REG), arg2(IMM|REG) : arg0 = arg1 << arg2 \\
Операция:  arg0 <- arg1 << arg2 \\
           FREG.DBIT <- вытесненный старший бит (бит 15) \\
           При SMOD=1 в младший бит подаётся предыдущий DBIT; при SMOD=0: 0 \\[2pt]
Пример:    BSL R0, R1, 1    ; Сдвиг R1 влево на 1 бит, результат в R0
}{Параметр сдвига: константа или динамическое значение регистра.}{Вытесненный старший бит фиксируется в FREG.DBIT (бит 13). При SMOD=1 в младший бит подаётся предыдущий DBIT; при SMOD=0 освободившиеся биты заполняются нулями.}{Элегантный кольцевой и потоковый сдвиг многобайтовых массивов.}
\InstructionPage{BSR}{11001}{Сдвиги}{Баррельный сдвиг вправо}{%
\small
Сигнатура: BSR arg0(REG), arg1(REG), arg2(IMM|REG) : arg0 = arg1 >> arg2 \\
Операция:  arg0 <- arg1 >> arg2 \\
           FREG.DBIT <- вытесненный младший бит (бит 0) \\
           При SMOD=1 в старший бит подаётся предыдущий DBIT; при SMOD=0: 0 \\[2pt]
Пример:    BSR R0, R1, 1    ; Сдвиг R1 вправо на 1 бит, результат в R0
}{Логический или арифметический сдвиг через выбор \field{cfg}.}{Младший вытесненный бит фиксируется в DBIT (бит 13). При SMOD=1 в старший бит вдвигается предыдущий DBIT; при SMOD=0 заполняется нулями.}{Потоковое масштабирование и непрерывная битовая фильтрация.}
\InstructionPage{MSP}{11010}{Система}{Манипуляция специальными регистрами (FREG, PTREG)}{%
\small
Сигнатура: MSP sreg, [op], args : работа со спец-регистрами \\
Операция:  
  \textbf{Прямая загрузка (бит 10 = 0):} \\
  \quad sreg=00 (FREG):  FREG <- Rs[15:0] (октет 1: [7:5]=Rs) \\
  \quad sreg=01 (PTREG): PTREG <- (Rs\_hi $\ll$ 16) | Rs\_lo (октет 1: [7:5]=Rs\_hi, [4:2]=Rs\_lo) \\
  \textbf{Модификация (бит 10 = 1):} \\
  \quad sreg=00 (FREG):  FREG <- FREG $\odot$ Rs (октет 1: [7:6]=op \{AND,OR,XOR,XNOR\}, [5:3]=Rs) \\
  \quad sreg=01 (PTREG): PTREG <- PTREG $\pm$ Rs (октет 1: [7]=op \{ADD,SUB\}, [6:4]=Rs) \\[2pt]
Пример:    MSP FREG, R0         ; Загрузка во FREG значения R0 \\
           MSP PTREG, R0, R1    ; Загрузка 24-битного PTREG = \{R0, R1\} \\
           MSP FREG, OR, R2     ; Наложение маски: FREG = FREG | R2 \\
           MSP PTREG, ADD, R3   ; Смещение адреса: PTREG = PTREG + R3
}{Регистры FREG (00) и 24-битный PTREG (01). Поддерживает прямую запись, логические маски FREG и локальное смещение PTREG.}{Флаги FREG модифицируются напрямую либо через маски. Арифметика PTREG формирует флаг переноса CF через АЛУ.}{Универсальный управляющий примитив системного уровня.}
\InstructionPage{RSV27}{11011}{Резерв}{Аппаратное расширение функционала}{%
\small
Сигнатура: RSV27 : зарезервировано \\
Операция:  Reserved for LIM architecture hardware extensions
}{Зарезервировано для будущих ревизий архитектуры LIM.}{Не определено.}{Зарезервировано.}
\InstructionPage{RSV28}{11100}{Резерв}{Аппаратное расширение функционала}{%
\small
Сигнатура: RSV28 : зарезервировано \\
Операция:  Reserved for LIM architecture hardware extensions
}{Зарезервировано для будущих ревизий архитектуры LIM.}{Не определено.}{Зарезервировано.}
\InstructionPage{RSV29}{11101}{Резерв}{Аппаратное расширение функционала}{%
\small
Сигнатура: RSV29 : зарезервировано \\
Операция:  Reserved for LIM architecture hardware extensions
}{Зарезервировано для будущих ревизий архитектуры LIM.}{Не определено.}{Зарезервировано.}
\InstructionPage{RSV30}{11110}{Резерв}{Аппаратное расширение функционала}{%
\small
Сигнатура: RSV30 : зарезервировано \\
Операция:  Reserved for LIM architecture hardware extensions
}{Зарезервировано для будущих ревизий архитектуры LIM.}{Не определено.}{Зарезервировано.}
\InstructionPage{RSV31}{11111}{Резерв}{Аппаратное расширение функционала}{%
\small
Сигнатура: RSV31 : зарезервировано \\
Операция:  Reserved for LIM architecture hardware extensions
}{Зарезервировано для будущих ревизий архитектуры LIM.}{Не определено.}{Зарезервировано.}

\clearpage
\section{Физический интерфейс и цоколёвка (34-выводный корпус)}
Процессор размещается в высоконадёжном компактном 34-выводном корпусе с двухрядным расположением выводов:

\begin{center}
\begin{tikzpicture}[font=\sffamily\scriptsize,
  pinbox/.style={draw,thick,fill=white,minimum width=6mm,minimum height=4mm,inner sep=0pt,font=\bfseries\tiny}]
\def\chipW{3.2}
\def\chipH{9.9}
\draw[thick,fill=limlight] (0,0) rectangle (\chipW,-\chipH);
\draw[thick] ($(\chipW/2,0)+(-0.35,0)$) arc (180:360:0.35);
\node[font=\bfseries\Large,text=limblue,rotate=-90] at (\chipW/2,-\chipH/2) {LIM M2};
\foreach \i/\sigL/\sigR in {
  1/A0/D0, 2/A1/D1, 3/A2/D2, 4/A3/D3, 5/A4/D4, 6/A5/D5, 7/A6/D6, 8/A7/D7,
  9/A8/IA, 10/A9/VCC, 11/A10/GND, 12/A11/RD, 13/A12/WD, 14/A13/RF, 15/A14/AL, 16/A15/RDI, 17/CLK1/CLK2
} {
  \pgfmathsetmacro{\y}{-\i * 0.55}
  \pgfmathtruncatemacro{\pinR}{35 - \i}
  \node[pinbox] at (0,\y) {\i};
  \node[anchor=east] at (-0.38,\y) {\sigL};
  \node[pinbox] at (\chipW,\y) {\pinR};
  \node[anchor=west] at (\chipW+0.38,\y) {\sigR};
}
\end{tikzpicture}
\end{center}

\subsection{Функциональное назначение внешних линий}
\begin{tabularx}{\textwidth}{>{\ttfamily}p{28mm}p{24mm}X}
\toprule
Вывод & Тип & Описание системного сигнала \\
\midrule
A0--A15    & Выход / Tri-State & 16-разрядная магистраль адреса (мультиплексируется с защёлкой AL). \\
D0--D7     & Двунаправленный   & 8-разрядная внешняя шина данных обмена с памятью и периферией. \\
CLK1, CLK2 & Вход              & Двухфазная тактовая синхронизация с защитой от перекрытия фаз. \\
VCC, GND   & Питание           & Основные шины питания логики (+5 В номинал). \\
RD, WD     & Выход             & Стробы чтения (Read Enable) и записи (Write Enable) шины. \\
AL         & Выход             & Address Latch: строб фиксации адреса во внешнем регистре. \\
RF         & Вход              & Ready Flag: аппаратная квитанция готовности ведомого устройства. \\
RDI        & Вход              & Raw Direct Interrupt: приоритетный вход аппаратного прерывания. \\
IA         & Выход             & Interruption Accepted / Interrupt Acknowledge: строб подтверждения прерывания устройствам; процессор ожидает считывания номера вектора с шины данных. \\
\bottomrule
\end{tabularx}

\clearpage
\section{Подсистема прерываний и векторная обработка}

Подсистема прерываний процессора LIM~M2 спроектирована для обеспечения минимальной латентности реакции на события в реальном времени, строгого детерминизма и полной защиты от шинных конфликтов. Архитектура объединяет аппаратные внешние запросы и программные исключения в единую 16-векторную иерархию с надёжным аппаратным квитированием на физическом уровне.

\subsection{Единое пространство 16 векторов прерываний}
Процессор поддерживает ровно 16 системных векторов (номера 0--15), которые являются ортогональными и общими как для внешней периферии, так и для внутренних событий ядра:

\begin{table}[h]
\centering
\caption{Единая таблица 16 векторов прерываний LIM M2}
\label{tab:vectors}
\begin{tabularx}{\textwidth}{>{\bfseries}c >{\ttfamily}c p{32mm}X}
\toprule
Вектор & Код & Источник / Тип & Функциональное назначение \\
\midrule
VEC 0  & 0x00 & Аппаратный / Сброс & Power-on Reset: начальная инициализация процессора и регистров. \\
VEC 1  & 0x01 & Аппаратный (NMI)   & Немаскируемое прерывание: сбой питания, аппаратная авария системы. \\
VEC 2  & 0x02 & Программная ловушка & Арифметическое исключение (деление на ноль, флаг \reg{FREG.EM}). \\
VEC 3  & 0x03 & Программная ловушка & Системный вызов ядра (System Call / Trap диспетчера ОС). \\
VEC 4  & 0x04 & Аппаратный / Периферия & Внешний таймер / Системный квантователь времени. \\
VEC 5  & 0x05 & Аппаратный / Периферия & Контроллер последовательного интерфейса (UART RX/TX). \\
VEC 6  & 0x06 & Аппаратный / Периферия & Контроллер прямого доступа к памяти (DMA Ready). \\
VEC 7  & 0x07 & Аппаратный / Периферия & Внешний контроллер прерываний / Шинный мост. \\
VEC 8--15 & 0x08--0x0F & Общие (HW / SW) & Пользовательские аппаратные линии и программные векторы. \\
\bottomrule
\end{tabularx}
\end{table}

\subsection{Аппаратный протокол квитирования (RDI $\to$ IA $\to$ D0--D7)}
Для взаимодействия с внешними контроллерами реализован надёжный 4-фазный цикл аппаратного рукопожатия (handshake), исключающий необходимость предварительного опроса регистров:

\begin{enumerate}
  \item \textbf{Фаза запроса (RDI):} Периферийное устройство выставляет высокий логический уровень на входе \texttt{RDI} (\textit{Raw Direct Interrupt}). Линия опрашивается ядром на границах микрокоманд.
  \item \textbf{Фаза проверки маски (Арбитраж):} Если глобальное маскирование прерываний разрешено (бит \reg{FREG.MIE} = 1), процессор фиксирует запрос и завершает текущую неделимую микрооперацию тракта данных.
  \item \textbf{Фаза подтверждения (Строб IA):} Ядро генерирует активный строб на выходе \texttt{IA} (\textit{Interruption Accepted / Interrupt Acknowledge}, вывод~26). Этот строб извещает все внешние устройства на шине о готовности ядра принять номер вектора.
  \item \textbf{Фаза передачи вектора (Шина D0--D7):} Одновременно с выдачей \texttt{IA} внешняя шина данных \texttt{D0--D7} переводится ядром в режим приёма (вход), а линия чтения \texttt{RD} активируется. Запросившее устройство удерживает на линиях \texttt{D0--D7} 8-битный идентификатор вектора (младшие 4 бита адресуют векторы 0--15).
  \item \textbf{Фаза фиксации и переход:} Процессор защёлкивает данные с шины \texttt{D0--D7}, снимает сигнал \texttt{IA} и переходит к процедуре сохранения контекста.
\end{enumerate}

\begin{center}
\begin{tikzpicture}[x=1.3cm,y=0.65cm,>=Stealth,font=\sffamily\scriptsize]
  \foreach \t in {0,...,6} {
    \draw[dotted,gray!50] (\t, 0.5) -- (\t, 5.5);
    \node[above,gray!70,font=\tiny] at (\t+0.5, 5.5) {T\t};
  }
  
  \node[left,font=\bfseries] at (0, 5) {CLK};
  \draw[thick,blue!80!black] (0, 4.7) 
    \foreach \t in {0,...,5} { -- ++(0.5, 0.6) -- ++(0.5, 0) -- ++(0, -0.6) -- ++(0, 0) };

  \node[left,font=\bfseries] at (0, 4) {RDI};
  \draw[thick,red!80!black] (0, 3.8) -- (1.0, 3.8) -- (1.2, 4.4) -- (4.8, 4.4) -- (5.0, 3.8) -- (6.0, 3.8);
  \node[above,red!80!black,font=\tiny] at (2.0, 4.4) {Запрос устройства};

  \node[left,font=\bfseries] at (0, 3) {IA};
  \draw[thick,teal!80!black] (0, 2.8) -- (2.0, 2.8) -- (2.2, 3.4) -- (4.2, 3.4) -- (4.4, 2.8) -- (6.0, 2.8);
  \node[above,teal!80!black,font=\tiny] at (3.2, 3.4) {Строб IA (Ack)};

  \node[left,font=\bfseries] at (0, 2) {D0--D7};
  \draw[thick,gray!60] (0, 1.8) -- (2.2, 1.8);
  \draw[thick,black] (2.2, 1.8) -- (2.4, 2.2) -- (4.0, 2.2) -- (4.2, 1.8) -- (6.0, 1.8);
  \draw[thick,black] (2.2, 1.8) -- (2.4, 1.4) -- (4.0, 1.4) -- (4.2, 1.8);
  \node[font=\bfseries\tiny] at (3.2, 1.8) {Вектор [3:0]};

  \node[left,font=\bfseries] at (0, 1) {ЯДРО};
  \draw[thick,cyan!80!black] (0, 0.8) -- (2.0, 0.8) -- (2.2, 1.2) -- (4.2, 1.2) -- (4.4, 0.8) -- (6.0, 0.8);
  \node[above,cyan!80!black,font=\tiny] at (1.0, 0.8) {Микрокод};
  \node[above,cyan!80!black,font=\tiny] at (3.2, 1.2) {Приём вектора};
  \node[above,cyan!80!black,font=\tiny] at (5.2, 0.8) {Стек (PC, FREG)};
\end{tikzpicture}
\end{center}

\subsection{Сохранение контекста в 24-битном стеке и вход в обработчик}
Принятие прерывания разрешено строго при условии, что процессор не находится в состоянии обработки другого прерывания (\reg{FREG.INTR} = 0) и прерывания глобально разрешены (\reg{FREG.MIE} = 1). Если \reg{FREG.INTR} = 1, любые повторные аппаратные прерывания аппаратно заблокированы.

При фиксации прерывания ядро выполняет следующую аппаратную процедуру:
\begin{enumerate}
  \item \textbf{Сохранение адреса возврата (PC):} Текущее значение счётчика команд \texttt{PC} сохраняется на вершине системного стека со смещением 24-битного указателя \reg{SREG}:
  $$\reg{SREG} \leftarrow \reg{SREG} - 2, \quad \text{Memory}[\reg{SREG}] \leftarrow \text{PC}_{15:0}$$
  \item \textbf{Сохранение флагов и статуса (FREG):} Текущее значение регистра \reg{FREG} (с исходным значением \reg{FREG.INTR}=0) помещается в стек:
  $$\reg{SREG} \leftarrow \reg{SREG} - 2, \quad \text{Memory}[\reg{SREG}] \leftarrow \reg{FREG}$$
  \item \textbf{Аппаратная блокировка вложенных прерываний:} Ядро выставляет флаг нахождения в прерывании \reg{FREG.INTR} $\leftarrow$ \texttt{1} (бит 15) и сбрасывает маску \reg{FREG.MIE} $\leftarrow$ \texttt{0}. Наличие активного бита \reg{FREG.INTR} гарантирует невозможность каскадных или паразитных повторных входов в обработчики.
  \item \textbf{Переход по вектору:} Базовый адрес обработчика вычисляется на основе принятого номера вектора и загружается в \texttt{PC}.
\end{enumerate}

\subsection{Возврат из прерывания}
Выход из процедуры обработки прерывания осуществляется восстановлением сохранённого контекста со стека:
\begin{enumerate}
  \item Со стека считывается и восстанавливается исходное состояние регистра флагов:
  $$\reg{FREG} \leftarrow \text{Memory}[\reg{SREG}], \quad \reg{SREG} \leftarrow \reg{SREG} + 2$$
  (При этом бит \reg{FREG.INTR} автоматически возвращается в \texttt{0}, снимая блокировку прерываний, а бит \reg{FREG.MIE} восстанавливает своё исходное значение).
  \item Со стека восстанавливается счётчик команд:
  $$\text{PC} \leftarrow \text{Memory}[\reg{SREG}], \quad \reg{SREG} \leftarrow \reg{SREG} + 2$$
  \item Процессор возобновляет выполнение прерванного потока инструкций со следующего такта без потерь промежуточных данных.
\end{enumerate}

\section{Заключение}
Архитектура \textbf{LIM M2} сочетает в себе проверенную надёжность многотактового микропрограммного управления и современную эффективность мультиплексорного тракта без третьего состояния. Такое решение обеспечивает чистейшие цифровые фронты, исключает паразитные токи и гарантирует непревзойдённую стабильность при масштабировании тактовых частот от 200 МГц до пределов физического кремния.
\end{document}
